\subsection{The precise parametric extent of the PL comparison}

\begin{saresult}{A.114}{Proved}{The concordance comparison is surjective on the fundamental group}
Let $B$ be a closed compact PL manifold, and use total-space PL and
definable families to define the respective simplicial concordance
groups $P_{\rm PL}(B)\subset P_{\rm def}(B)$. In addition to the
bijection on components in A.113, inclusion induces a surjection
\[
 \pi_1(P_{\rm PL}(B),\operatorname{id})
 \longrightarrow
 \pi_1(P_{\rm def}(B),\operatorname{id}).
\]
More explicitly, every based definable circle family of concordances
is based definably homotopic to a family that is PL as a map of
total spaces.
If the field is $\R$ and every component of $B$ has dimension at
least five, this map on $\pi_1$ is an isomorphism. In that scope
the inclusion is injective on every homotopy group.
For $k\ge2$ there is the precise splitting
\[
 \pi_kP_{\rm def}(B)\simeq
 \pi_kP_{\rm PL}(B)\oplus
 \ker\bigl(\pi_kP_{\rm def}(B)\to\pi_kP_{\rm TOP}(B)\bigr).
\]
\end{saresult}
\begin{proof}
Write $Z=B\times[0,1]$ and let $C$ be a finite polygonal circle
with a distinguished vertex $c_0$. A loop is a homeomorphism
\[
 H:Z\times C\longrightarrow Z\times C
\]
over $C$, fixing $B\times\{0\}\times C$ pointwise and equal to
the identity over $c_0$. Mark also the upper boundary component.
Choose a PL function $a:C\to K$ nonconstant on every edge; assigning
distinct values to the finitely many vertices suffices. Apply the
complement to Shiota's Theorem 2.2 to the scalar function
$f=a\circ\operatorname{pr}_C$, with its two triangulations
$\pi=\operatorname{id}$ and $\pi'=H$. Mark the compact subpolyhedra
$Z\times e$ for every closed edge $e$ of $C$, and the two boundary
components of $Z$ times $C$. The source statement is
\cite[Theorem 2.2, Complement, p.~7]{Shiota2010}.

It supplies a definable isotopy $\omega_s$ from $H$ to a PL
homeomorphism, preserving $f$ and every marked set. Each $\omega_s$
preserves the full circle coordinate: for a point with coordinate
$c\in e$, its image has coordinate in $e$ with the same value of
$a$, and $a|_e$ is injective. This also works at vertices.
Thus $\omega_s$ is a family of homeomorphisms of $Z$ over $C$.
It need not yet fix the bottom pointwise. Write its bottom
restriction as
\[
 (b,0,c)\longmapsto(\alpha_s(b,c),0,c)
\]
and extend this restriction over the interval by
\[
 E_s(b,r,c)=(\alpha_s(b,c),r,c).
\]
This is a definable family of homeomorphisms. At $s=1$ it is
PL, since it is the product extension of the restriction of the
PL map $\omega_1$. The homeomorphisms
$g_s=E_s^{-1}\omega_s$ therefore fix the bottom and preserve the
top. At $s=0$ the correction is the identity, so $g_0=H$.
At $s=1$, $g_1$ is PL as a total-space map.

Finally let $q_s=g_s|_{Z\times\{c_0\}}$, viewed as a concordance
of $B$, and replace $g_s$ by
\[
 h_s=(q_s^{-1}\times\operatorname{id}_C)\circ g_s.
\]
Then $h_s$ is based at the identity
over $c_0$ for every $s$. Since $q_0=\operatorname{id}$, it starts
at $H$. Since $q_1$ is a restriction of a PL map and its product
extension is PL, it ends at a jointly PL family. These constructions
are continuous on total spaces and definable; inverse families
are continuous by compactness, or directly by the displayed
compositions. They prove the asserted based loop homotopy.

Loops and homotopies in these simplicial groups are represented,
after finite simplicial subdivision, by precisely such finite
parameter families. Thus the statement about families proves
surjectivity on $\pi_1$ of their realizations.

Over $\R$ a definable family is a continuous family and therefore
defines a map to $P_{\rm TOP}(B)$. The composite
\[
 P_{\rm PL}(B)\longrightarrow P_{\rm def}(B)
 \longrightarrow P_{\rm TOP}(B)
\]
is the usual PL--TOP comparison. Under the dimension hypothesis
it is a weak equivalence by A.113. Its maps on homotopy groups
are injective, so the first factor is injective on every homotopy
group. The surjectivity already proved gives the assertion on
$\pi_1$. This argument is confined to the real field, where
the displayed forgetful map to topological families exists.
If $i_k$ and $j_k$ denote the two maps on $\pi_k$ in that
display, $(j_ki_k)^{-1}j_k$ is a retraction of $i_k$.
For $k\ge2$ these groups are abelian, so its kernel gives the
stated direct-sum splitting. It is exactly $\ker j_k$.
\end{proof}

\begin{saresult}{A.115}{Proved}{All PL concordance homotopy groups are independent of the real closed field}
Let $F\subset K$ be an extension of real closed fields and let $B$
be a closed compact PL manifold given by a finite polyhedral model
over $F$. Extension of coefficients induces a weak equivalence
of simplicial groups
\[
 P_{{\rm PL},F}(B)\longrightarrow P_{{\rm PL},K}(B_K).
\]
Consequently a fixed finite abstract PL model has the same PL
concordance homotopy type over any two real closed fields, by the
zigzag through a realization over the field of real algebraic
numbers. This is a zigzag of coefficient-extension maps; it is
not a literal map from a non-Archimedean field's families to
ordinary topological families.
\end{saresult}
\begin{proof}
We give the parameter-preserving argument. A statement only about
unparametrized homeomorphisms would not suffice.

\emph{Extension on product cells.}
Let $Z$ be the underlying space of a finite complex of convex
cells, let $Q$ be a finite parameter polyhedron and $Q_0$ a
subcomplex. Suppose a PL homeomorphism of $Z\times Q_0$ is over
$Q_0$ and preserves each cell of $Z$. It extends to a PL
homeomorphism of $Z\times Q$ over $Q$, still preserving every
cell. More generally it extends relative to any already
specified union of product cells on which the same conditions hold.
To see this, triangulate $Q$ and proceed over the cells
$\sigma\times\tau$ in increasing total dimension. On the boundary
of a cell not already specified the homeomorphism has been
defined, maps that boundary to itself, and preserves the parameter
coordinate. Choose an interior point
$c=(c_\sigma,c_\tau)$ of this convex product cell and extend
by coning:
\[
 c+t(z-c)\longmapsto c+t(H(z)-c),
 \qquad z\in\partial(\sigma\times\tau),\quad 0\le t\le1.
\]
This is PL after starring a triangulation of the boundary at
$c$, and is a homeomorphism by the same construction for the
inverse boundary map. Its parameter coordinate is unchanged,
since
\[
 \operatorname{pr}_Q(c+t(H(z)-c))
 =(1-t)c_\tau+t\operatorname{pr}_Q(z)
 =\operatorname{pr}_Q(c+t(z-c)).
\]
Thus it is coning on the \emph{total product cell}, not separate
radial coning in each fibre. The maps agree on faces and the
finite induction proves the claim. It works over every ordered
field. In particular cell-preserving PL families can be contracted
through such families. Two families agreeing on a union of
product cells can be contracted with that agreement retained:
first perform this same extension on the common part, then on
the remaining cells separately.

Put $Z=B\times[0,1]$, with a fixed finite convex-cell
decomposition $\mathcal C_Z$ over $F$ that contains its bottom
and top as subcomplexes. Represent a homotopy group element by
a PL homeomorphism
\[
 h:Z_K\times P_K\longrightarrow Z_K\times P_K,
 \qquad P=[0,1]^q,
\]
over $P_K$, fixing the bottom and equal to the identity on
$Z_K\times\partial P_K$. For $q=0$ the last condition is empty.
Fix a cell decomposition $\mathcal C_P$ of $P$ over $F$.

\emph{Finite transferred data.}
Apply simultaneous subdivision to the one-way chain consisting
of two distinct copies of the total space,
\[
 Z_K\times P_K\xrightarrow{h}Z_K\times P_K
          \xrightarrow{\operatorname{pr}_P}P_K.
\]
The source and target subdivisions $M,N$ can be required to
refine the same fixed product cells
$\mathcal C_Z\times\mathcal C_P$, and the base subdivision $L$
refines $\mathcal C_P$. Then $h:M\to N$ is a simplicial
isomorphism and both projections to $L$ are simplicial, the
source projection being the composite of the displayed maps.
This is exactly the permitted one-way case of
\cite[Theorems 2.14--2.15, pp.~17--19]{RourkeSanderson1972};
it does not ask for a simplicial subdivision of a cycle.
On the bottom and on $Z_K\times\partial P_K$, where $h$ is
the identity, the induced source and target subdivisions coincide.

Keep the finite abstract complexes and all their carrier labels
fixed, but regard their vertex coordinates as unknowns. The
requirements that they realize subdivisions of the designated
fixed product cells, that the abstract map $M\to N$ is realized
by a PL homeomorphism, and that each source vertex and its
corresponding target vertex project to their designated vertex
of $L$, are a finite first-order condition over $F$. To spell
out the finiteness, affine independence is expressed by nonzero
minors; simplex membership uses finitely many barycentric
coordinates; intersection along the prescribed common faces
and coverage of each fixed cell are finite quantified field
formulas. Add equality of the corresponding source and target
vertices on the bottom and on the parameter boundary. The
given configuration satisfies this condition over $K$.
Quantifier elimination for real closed fields, hence
$F\preccurlyeq K$, gives a configuration $M',N',L'$ over $F$
satisfying the same condition. It determines a PL family
$h_F$ of concordances, based on $\partial P$, over $F$.

There are matching-vertex PL homeomorphisms over $K$
\[
 u:M\longrightarrow M',\qquad
 v:N\longrightarrow N',\qquad
 \beta:L\longrightarrow L',
 \qquad h_F=vhu^{-1}.
\]
Each of $u,v$ preserves every fixed product cell, and $\beta$
preserves every fixed base cell. The projection identities
hold on entire simplices, not just their vertices:
\[
 \operatorname{pr}_P u=\beta\operatorname{pr}_P,\qquad
 \operatorname{pr}_P v=\beta\operatorname{pr}_P.
\]
Indeed a simplex of $M$ projects into one simplex of $L$,
on which $\beta$ is affine. For
$x=\sum_i\lambda_i x_i$ in this simplex, the vertex-incidence
equalities give
$\operatorname{pr}_P u(x)=\sum_i\lambda_i\beta(\operatorname{pr}_P x_i)
=\beta(\operatorname{pr}_P x)$; the proof for $N$ is identical.

Set
\[
 u_0=(\operatorname{id}_Z\times\beta^{-1})u,\qquad
 v_0=(\operatorname{id}_Z\times\beta^{-1})v.
\]
They preserve the parameter pointwise and every product cell.
They agree on the bottom and on $Z_K\times\partial P_K$,
because the matching-vertex maps agree there. Apply the
product-cell extension construction over the enlarged parameter
$P_K\times[0,1]$ to obtain PL isotopies $u_s,v_s$ from the
identity to $u_0,v_0$ which retain this agreement for every $s$.
Concretely extend their common boundary data first on the union
of these product subcomplexes, then extend over the other cells.
The families
\[
 v_s h u_s^{-1}
\]
are therefore concordances fixed on the bottom and based on
the parameter boundary throughout. Their last member is
$(\operatorname{id}_Z\times\beta^{-1})h_F
(\operatorname{id}_Z\times\beta)$.
The same product-cell construction gives an isotopy of the
cell-preserving map $\beta$ to the identity, preserving every
base cell. Conjugating by that isotopy connects this last member
to $h_F$ while keeping the parameter boundary based. This proves
surjectivity on every homotopy group, including components.

\emph{Injectivity.}
Suppose a based PL family over $F$ becomes nullhomotopic over
$K$. Take one finite PL triangulation of its $K$-nullhomotopy.
Its vertex coordinates again give a finite first-order
configuration problem over $F$, with the extra requirement that
on the parameter boundary its PL map equal the prescribed
$F$-family or the identity, as appropriate. That equality is
itself a finite field formula, because the prescribed family
has finitely many affine pieces over $F$. Transfer therefore
gives an $F$-PL nullhomotopy with exactly the same boundary data.
The same argument applies to a path between two zero-simplices.
Thus the maps on all homotopy groups are injective.

Finally simplicial groups are Kan, and their homotopy groups
are computed by these finite parameter families. The established
isomorphisms prove the weak equivalence. Real algebraic numbers
embed as an ordered real closed subfield in every real closed
field; realizing a finite abstract complex by standard rational
vertices gives the asserted comparison zigzag.
\end{proof}

\begin{saresult}{A.116}{Proved}{Semialgebraic concordance comparison reduces to the real field}
Let $F\subset K$ be real closed fields and let $B$ be a closed
compact PL manifold given by a finite polyhedral model over $F$.
Extension of coefficients induces a weak equivalence
\[
 P_{{\rm SA},F}(B)\longrightarrow P_{{\rm SA},K}(B_K),
\]
where families and their homotopies are semialgebraic maps of
total spaces. The squares formed by these maps and the PL
coefficient-extension maps commute. Thus the homotopy-group
comparison $P_{\rm PL}\to P_{\rm SA}$ for a finite abstract PL
model is independent of the real closed field and reduces to $\R$.

In particular, if every component of $B$ has dimension at least
five, then over every real closed field this comparison is
bijective on components, is an isomorphism on $\pi_1$, and is
injective on every homotopy group. Full weak equivalence is not
asserted here.
\end{saresult}
\begin{proof}
We give the continuity argument needed in addition to transfer.
Merely transferring the existence of a homeomorphism would not
connect its coefficient tuple to the original family.

\emph{Continuous evaluation after a finite coefficient partition.}
Let $D$ be a closed bounded semialgebraic set over $F$, and fix
one ordered-field formula $\varphi(a;x,y)$ over $F$. Let $T$
consist of the tuples $a$ for which this formula is the graph of
a homeomorphism $h_a:D\to D$ satisfying a prescribed finite
list of first-order conditions. Totality, uniqueness,
bijectivity, and continuity of both the map and its inverse are
first-order conditions: continuity uses quantifiers over points
and positive $\epsilon,\delta$ in the field. Quantifier
elimination therefore makes $T$ semialgebraic over $F$.

Form the graph space
\[
 \Gamma=\{(a,x,y):a\in T,\ x,y\in D,\ \varphi(a;x,y)\}.
\]
Its projection to $T$ is continuous and semialgebraic, although
evaluation on $T\times D$ need not yet be continuous. Apply
semialgebraic triviality to this projection over $F$, and refine
the resulting base partition into semialgebraic cells over $F$.
The precise input is Hardt triviality without a properness
assumption on the original map; its real-closed-field form is
stated in \cite[Section 2.2.1, p.~1296]{LerarioRizziTiberio2025}.
For each cell $C$, choose $a_0\in C(F)$ and identify the fixed
fibre of its trivialization with the fibre at $a_0$. We obtain
\[
 \theta:C\times\Gamma_{a_0}\xrightarrow{\ \sim\ }\Gamma|_C
\]
over $C$. Choose all these data over $F$ before extending to
$K$. Their finite first-order properties, including continuity,
remain true after coefficient extension.

The fibre $\Gamma_{a_0}$ is closed and bounded, being the graph
of a continuous map on $D$. Compose $\theta$ with the projection
which forgets the output coordinate to obtain a continuous
semialgebraic bijection
\[
 J:C\times\Gamma_{a_0}\longrightarrow C\times D.
\]
It is a homeomorphism. Indeed, for a closed bounded
semialgebraic $A\subset C\times D$, its projection
$C_0\subset C$ is closed and bounded, and $J^{-1}(A)$ is
closed in $C_0\times\Gamma_{a_0}$, hence closed and bounded.
Thus $J$ is definably proper. Inverse continuity can also be
checked directly: curve selection applied to a failure of
continuity would give a semialgebraic curve tending to a point
of $C\times D$ whose inverse fails to tend to the unique
inverse point. Its second inverse coordinate lies in
$\Gamma_{a_0}$ and has a limit by boundedness and semialgebraic
monotonicity. Continuity of $J$ identifies that limit with the
unique inverse point, a contradiction. This uses definable
compactness, not ordinary compactness of the order topology
of a non-Archimedean field.

The output projection composed with $\theta J^{-1}$ is
$(a,x)\mapsto h_a(x)$. Thus evaluation is jointly continuous on
$C\times D$. Interchanging the input and output coordinates
proves the same assertion for inverse evaluation. We have
obtained genuine families of homeomorphisms on the coefficient
cells.

\emph{Surjectivity on homotopy groups.}
Put $Z=B\times[0,1]$ and $D=Z\times P$, where $P=[0,1]^q$.
Represent a homotopy-group class over $K$ by a semialgebraic
homeomorphism $h:D_K\to D_K$ over $P_K$, fixed on the bottom
and equal to the identity over $\partial P_K$. For $q=0$ the
boundary condition is empty. Choose one finite formula for its
graph and replace its finitely many extra coefficients by a
tuple $a$, leaving $D$ and the formula fixed over $F$. In $T$
impose exactly the conditions just stated: preservation of each
parameter point, the bottom fixed pointwise, the top preserved,
and the identity on the parameter boundary.

The original tuple lies in $T(K)$, hence in $C(K)$ for one of
the cells chosen over $F$. Choose $a_0\in C(F)$. A semialgebraic
cell has a semialgebraic homeomorphism over $F$ with an open
cube (or is a point). Transporting a line segment in that cube
gives a semialgebraic path in $C(K)$ from the original tuple to
$a_0$. Pull back the jointly continuous evaluations along this
path. They give a semialgebraic homotopy of concordance families
through members satisfying all conditions encoded in $T$.
The last family is defined over $F$. This proves surjectivity
on every homotopy group and on components.

\emph{Injectivity on homotopy groups.}
Suppose a based family over $F$ admits a semialgebraic
nullhomotopy over $K$. Choose one finite graph formula for that
nullhomotopy and replace its extra coefficients by variables.
Require a total continuous family of homeomorphisms with
continuous inverse, preserved parameters, fixed bottom and
parameter sides, and exact equality with the prescribed
$F$-family and the identity at the two endpoints. These are
first-order conditions over $F$, including equality with the
prescribed endpoint graph. A solution exists in $K$, so
$F\preccurlyeq K$ supplies a solution in $F$. This is the
required $F$-semialgebraic nullhomotopy, with its boundary
unchanged. The same argument descends paths between two
zero-simplices. Thus every homotopy-group map is injective.
As in A.115, the Kan property and finite parameter
representatives give weak equivalence of simplicial groups.

Coefficient extension preserves PL maps, so the PL/SA squares
commute literally. Realize the finite abstract PL model over
the real algebraic numbers and use their embeddings in $K$
and $\R$. A.115 and the present theorem give vertical weak
equivalences in the comparison zigzag. Over $\R$, A.114 proves
injectivity in every degree and surjectivity on $\pi_1$ in the
stated dimension range; A.113 supplies the component assertion.
These assertions therefore hold over $K$. This reduction is
for semialgebraic structures; it does not assume that an
arbitrary o-minimal expansion has a model whose field is $\R$.
\end{proof}

\paragraph{Exact remaining gate; no full parametric conclusion.}
The scalar straightening proof in A.114 uses the fact that a scalar function is injective on
each edge of a finite graph. It does not generalize in this way
to a two-dimensional parameter disk. For an arbitrary o-minimal
expansion of a real closed field it does not prove injectivity on $\pi_1$, and it does not prove that
$P_{\rm PL}(B)\to P_{\rm def}(B)$ is a weak equivalence.
A first missing assertion is this relative extension statement:
given a definable family on $D^2$ whose boundary family is jointly
PL, can it be deformed relative to its boundary to a jointly PL
family while preserving every point of $D^2$ as a parameter?
Over $\R$ in the dimension range of A.113, the additional TOP
argument above already settles injectivity. By A.115--A.116 this
also holds for semialgebraic families over every real closed field.
In this semialgebraic scope the remaining comparison reduces to
the real field, and the first remaining group assertion is
surjectivity on $\pi_2$. This requires
straightening a sphere family with two independent parameters,
or another argument that rules out a definable class killed by
the forgetful map to TOP.

Triangulating a total-space homeomorphism with the finitely many
parameter faces marked only preserves those faces setwise. It
does not preserve every fibre of the projection to $D^2$.
Successively applying the scalar theorem to its two coordinates
does not repair the problem, because the second isotopy need not
preserve the first coordinate. General homotopy comparison for
maps of definable sets likewise does not keep intermediate maps
homeomorphisms. These are exact failures of proposed deductions,
not counterexamples to the desired weak equivalence. The obstruction
to triangulating arbitrary maps with a two-dimensional target is
explicitly discussed in \cite[p.~1231]{FichouShiota2017}; the weak
triangulation used there is local and does not supply the missing
parameter-preserving family theorem.

The local-contractibility arguments used for A.117 and A.209 take
place in TOP homeomorphism groups. Their continuous contractions
and quotient lifts are not asserted semialgebraic. In particular,
using such a contraction to fill a small semialgebraic sphere would
require a separate proof that the filling is jointly semialgebraic.
Individual Hauptvermutung, cellwise Hardt triviality, or a collection
of unrelated PL approximations does not establish this. The
compact-bundle reduction suggested by clutching a sphere family
has the same missing parameter-preserving reduction, so it is not
an independent proof.

The paragraph after Theorem 2.6 in
\cite[p.~35]{Shiota2010} proposes a PL reduction for definable
fibre bundles as a further possibility. Its Remark 2.7 proves
only stable equivalence for microbundles. Neither statement is
the missing relative two-parameter theorem. The sources opened
for this check, together with the arguments above, justify
A.113--A.116, but do not
settle the full higher-homotopy comparison.

\subsection{Pointwise cone families with locally parameterized outer faces}

There is a useful intermediate family convention between a derived
choice space and Hausdorff continuity of subsets alone. Here an outer
face is an unparameterized subset. A finite-parameter family is required
locally on its parameter space to admit jointly continuous embeddings
of a fixed abstract outer face onto those subsets. The parameterizations
are witnesses of continuity, not additional data retained in the objects.
In a definable simplicial family model the witnesses on each simplex
are jointly definable; finite nullhomotopies use finitely many such
simplex presentations. No finite subcover assertion over a
non-Archimedean field is implicit in this convention.

\begin{saresult}{A.117}{Refuted}{The cone obstruction persists for locally parameterized outer-face families}
Let $B$ be a closed connected PL manifold of dimension at least four,
let $X=C_2B$ be its cone with radial interval $[0,2]$, and put $C=C_1B$.
Give the locally flat outer faces of strict cone neighborhoods the
local quotient topology from
\[
 \operatorname{Emb}_{\rm lf}\bigl(B,B\times(0,2)\bigr)
       /\Homeo(B),
\]
where embeddings have the compact-open topology and $\Homeo(B)$ acts
by precomposition. An outer face determines the side containing the
vertex. The component $\mathcal W$ of $C$ has weak homotopy type
\[
             \mathcal W\simeq BP_{\rm TOP}(B).
\]
For $B=(S^1)^6$, it is therefore not contractible. Its nontrivial loop
can be chosen semialgebraic, and it cannot be filled by real definable
families of strict cones whose outer faces admit local joint
parameterizations. The same non-filling conclusion holds for finite
semialgebraic simplex families with these joint semialgebraic witnesses
over every real closed field.

If an external parameterization $B\to\dot C$ is retained as part of
each object, the component instead has weak homotopy type
\[
 B\Homeo\bigl(B\times[1,2]\ \mathrm{rel}\
               B\times\{1,2\}\bigr).
\]
This is a different convention from the unparameterized outer faces
in the preceding assertion.
\end{saresult}
\begin{proof}
We first record the topological lifting fact used below. If $G$ is a
metrizable topological group and $J$ is a closed, completely metrizable
subgroup that is locally $n$-connected for every finite $n$, every
continuous map $[0,1]^m\to G/J$ lifts to $G$. This is the finite-cube
lifting argument in the proof of
\cite[Corollary~2.2, pp.~4--5]{BanakhRepovs2010}. The additional
``balanced'' hypothesis in that corollary is used for its separate
homogeneous-metric conclusion, not for this lifting argument. Its
hypotheses can be checked directly: the cosets are closed complete
subsets in the two-sided completion, and the coset family is equi-locally
$n$-connected. For the latter, given an identity neighborhood $U$, choose
$W$ so that every sphere in $J\cap W$ contracts in $J\cap U$, and then
choose $V$ with $V^{-1}V\subset W$. Translating a sphere contained in
$gJ\cap xV$ by the inverse of any one of its points reduces it to that
contraction in $J$. The finite-dimensional selection theorem used in
the cited proof supplies lifts over cubes. In particular the assertion
uses complete metrizability, not completeness of a one-sided invariant
metric. A lift of a homotopy of cubes can be made to extend a prescribed
initial lift: multiply it on the right by its initial discrepancy in
$J$, constant in the homotopy coordinate. Thus $G\to G/J$ is a Serre
fibration.

Now set
\[
 G=\Homeo(X\ \mathrm{rel}\ \{o\}\cup B\times\{2\}),
 \qquad J=\{g\in G:g(C)=C\}.
\]
The Alexander contraction is a continuous strong contraction of $G$.
Restriction to the shell $B\times[1,2]$ gives a split map
$J\to P_{\rm TOP}(B)$, with section obtained by coning the inner-face
restriction. Its kernel is
\[
 K=\Homeo(C\ \mathrm{rel}\ \{o\}\cup B\times\{1\}).
\]
Restriction, conical extension, and composition are continuous in the
compact-open topology, including inverses. Hence this splitting is a
homeomorphism of underlying spaces $J\cong K\times P_{\rm TOP}(B)$.
The factor $K$ has the continuous strong Alexander contraction.
Such a contraction proves local contractibility at the identity:
for a specified identity neighborhood, continuity and compactness of
the contraction interval provide a smaller neighborhood whose entire
contraction stays in the specified one. Translation gives the assertion
at every point.

The other factor is locally contractible as well. The boundary-component
preserving subgroup of $\Homeo(B\times[1,2])$ is open and is locally
contractible by \cite[Corollary~1.1, p.~63]{EdwardsKirby1971}.
Restriction to the top boundary has the continuous section given by
constant extension across the interval. The resulting continuous
retraction onto the top-fixed subgroup shows that $P_{\rm TOP}(B)$
is locally contractible. Finally the metric
\[
 d_G(g,h)=\sup_{x\in X}d(gx,hx)
           +\sup_{x\in X}d(g^{-1}x,h^{-1}x)
\]
is complete for any complete compatible metric on the compact $X$.
A Cauchy sequence has mutually inverse uniform limits. The prescribed
fixed-point conditions and the equation $g(C)=C$ are closed in this
metric. Thus $G$ and $J$ meet the lifting hypotheses.

We next verify that the orbit topology is the stipulated outer-face
topology. The codimension-one local unknotting theorem
\cite[Theorem~7.3.1, p.~358]{DavermanVenema2009} says that sufficiently
close locally flat embeddings of the compact face are related by an
arbitrarily small ambient push. Here closeness is uniform closeness of
\emph{parameterized embeddings}. An $\epsilon$-push is supported in
the $\epsilon$-neighborhood of the face
\cite[p.~115]{DavermanVenema2009}. Choose $\epsilon$ smaller than its
distance from both $o$ and $B\times\{2\}$. The push then extends by the
identity to $X$ and belongs to $G$; its inverse is equally small. It
sends the cone side to the cone side because it fixes $o$.

More explicitly, let $U\subset G$ be open and $g\in U$. For every
parameterization $e:B\to g(\dot C)$, a sufficiently close embedding
$e'$ is $he$ for a push $h$ sufficiently close to the identity that
$hg\in U$. Therefore the inverse image of the orbit image of $U$ in
the parameterized embedding space is open. Passing to the quotient by
precomposition proves that $G\to G\cdot C$ is open. It consequently
induces a homeomorphism $G/J\cong G\cdot C$. The same argument at any
strict cone with a locally flat outer face shows that every ambient
orbit is open. As $G$ is contractible, its orbit is connected and is
exactly the component $\mathcal W$ under discussion.

The Serre fibration $J\to G\to\mathcal W$, with contractible middle
space, gives the asserted weak homotopy type and, in particular,
\[
 \pi_{k+1}(\mathcal W,C)\cong\pi_kJ
          \cong\pi_kP_{\rm TOP}(B)\qquad(k\geq0).
\]
Equivalently, the map from the homotopy quotient of the contractible
$G$ by $J$ to $G/J$ is a weak equivalence, identifying it with $BJ$.
No local normal charts preserving all family parameters were assumed.

For the explicit obstruction take a PL concordance $p$ of the six-torus
representing a nonzero class in $\pi_0P_{\rm TOP}(B)$; such a class and
its PL representative are supplied by A.103 and A.113. Regard $p$ as
an automorphism of the shell fixed on its top, and cone its inner-face
map to obtain $j\in J$. Let $H_s$, $0\leq s\leq1$, be the Alexander
isotopy from the identity to $j$. Then
\[
                  C_s=H_s(C)
\]
is a semialgebraic loop of strict cones, since $C_0=C_1=C$.
It has jointly continuous face parameterizations on the parameter
interval; the two endpoint parameterizations need not coincide,
because they are not part of the unparameterized objects.

Suppose this loop bounded a disk of the indicated real definable
families. Each individual definable outer face is TOP locally flat:
compatible triangulation makes the ambient annulus and the face
polyhedra, the definable Hauptvermutung makes both PL manifolds, and
a codimension-one PL embedding is TOP locally flat. The last statement
follows by induction on links and the generalized Schoenflies theorem,
as in \cite[Theorem~1.2.1, p.~29]{DavermanVenema2009}.
Local joint face parameterizations make the disk map continuous in the
embedding-quotient topology. Its image lies in $\mathcal W$, because
ambient orbits are open. By the finite-cube lifting fact the disk map
has a continuous lift $L$ to $G$. On its boundary compare $L_s$ with
$H_s$. The path $a_s=L_s^{-1}H_s$ lies in $J$, and $L_0=L_1\in J$.
Thus $a_s$ joins $L_0^{-1}$ to $L_0^{-1}j$, so multiplication by $L_0$
gives a path from the identity to $j$ in $J$. Shell restriction would
make $p$ trivial in $\pi_0P_{\rm TOP}(B)$, a contradiction. A disk
assembled from finitely many such parameterized simplices gives the
same continuous disk map, so the argument also excludes nullhomotopies
in the associated simplicial family model.

For semialgebraic families over an arbitrary real closed field, choose
the torus and $p$ over the real algebraic numbers, using A.115 and the
nonzero real PL class just used. If a finite semialgebraic filling
existed over another real closed field, write its finitely many graph
formulas and coefficient tuples. Joint continuity, fibrewise embedding
of each parameterized face, avoidance of the vertex and outer base,
agreement of face images on shared parameter faces, and the prescribed
boundary images are all first-order conditions. Transfer gives the
same finite face family over the real algebraic numbers and then over
$\mathbb R$. The transferred individual faces are semialgebraic and
therefore TOP locally flat. Conicality of the transferred interiors is
not needed: all these faces lie in the same open ambient orbit as the
boundary loop. This contradicts the real argument. No assertion about
real models of arbitrary o-minimal expansions is used.

Finally, retaining an external parameterization replaces $J$ by the
subgroup $J^{\rm pt}$ fixing $\dot C$ pointwise. Its shell restriction
has kernel $K$ and target
$Q=\Homeo(B\times[1,2]\ \mathrm{rel}\ B\times\{1,2\})$.
The group $Q$ is locally contractible: double the shell along its
boundary, apply \cite[Corollary~7.3, p.~83]{EdwardsKirby1971} relative
to the two doubling faces, and restrict the resulting open
component-preserving subgroup, which is $Q\times Q$, to one half.
All metric and orbit arguments above apply with $J^{\rm pt}$ and
parameterized embeddings. The component is therefore weakly equivalent
to $BQ$. This verifies exactly why that additional convention changes
the stabilizer.
\end{proof}

\paragraph{The distinct Hausdorff-only convention.}
A.117 settles the comparison for outer-face families locally continuous
in the compact-open embedding quotient. A.118 supplies an explicit
semialgebraic family of polygonal disks proving that Hausdorff
continuity of both cones and outer faces does not imply local joint
face parameterizations. The two conventions are therefore genuinely
different. Whether the six-torus loop becomes fillable in the weaker
Hausdorff-only space is a separate comparison question; it is not
needed for the refutation of II b) under the simplicial-category
convention stipulated in the prompt, already proved in A.111.
Likewise A.117 supplies TOP lifts, not definable lifts, and therefore
does not settle the higher PL-to-definable concordance comparison in
A.114--A.116.
